\IfFileExists{stacks-project.cls}{%
\documentclass{stacks-project}
}{%
\documentclass{amsart}
}

% For dealing with references we use the comment environment
\usepackage{verbatim}
\newenvironment{reference}{\comment}{\endcomment}
%\newenvironment{reference}{}{}
\newenvironment{slogan}{\comment}{\endcomment}
\newenvironment{history}{\comment}{\endcomment}

% For commutative diagrams we use Xy-pic
\usepackage[all]{xy}

% We use 2cell for 2-commutative diagrams.
\xyoption{2cell}
\UseAllTwocells

% We use multicol for the list of chapters between chapters
\usepackage{multicol}

% This is generally recommended for better output
\usepackage{lmodern}
\usepackage[T1]{fontenc}

% For cross-file-references
\usepackage{xr-hyper}

% Package for hypertext links:
\usepackage{hyperref}

% For any local file, say "hello.tex" you want to link to please
% use \externaldocument[hello-]{hello}
\externaldocument[introduction-]{introduction}
\externaldocument[conventions-]{conventions}
\externaldocument[sets-]{sets}
\externaldocument[categories-]{categories}
\externaldocument[topology-]{topology}
\externaldocument[sheaves-]{sheaves}
\externaldocument[sites-]{sites}
\externaldocument[stacks-]{stacks}
\externaldocument[fields-]{fields}
\externaldocument[algebra-]{algebra}
\externaldocument[brauer-]{brauer}
\externaldocument[homology-]{homology}
\externaldocument[derived-]{derived}
\externaldocument[simplicial-]{simplicial}
\externaldocument[more-algebra-]{more-algebra}
\externaldocument[smoothing-]{smoothing}
\externaldocument[modules-]{modules}
\externaldocument[sites-modules-]{sites-modules}
\externaldocument[injectives-]{injectives}
\externaldocument[cohomology-]{cohomology}
\externaldocument[sites-cohomology-]{sites-cohomology}
\externaldocument[dga-]{dga}
\externaldocument[dpa-]{dpa}
\externaldocument[sdga-]{sdga}
\externaldocument[hypercovering-]{hypercovering}
\externaldocument[schemes-]{schemes}
\externaldocument[constructions-]{constructions}
\externaldocument[properties-]{properties}
\externaldocument[morphisms-]{morphisms}
\externaldocument[coherent-]{coherent}
\externaldocument[divisors-]{divisors}
\externaldocument[limits-]{limits}
\externaldocument[varieties-]{varieties}
\externaldocument[topologies-]{topologies}
\externaldocument[descent-]{descent}
\externaldocument[perfect-]{perfect}
\externaldocument[more-morphisms-]{more-morphisms}
\externaldocument[flat-]{flat}
\externaldocument[groupoids-]{groupoids}
\externaldocument[more-groupoids-]{more-groupoids}
\externaldocument[etale-]{etale}
\externaldocument[chow-]{chow}
\externaldocument[intersection-]{intersection}
\externaldocument[pic-]{pic}
\externaldocument[weil-]{weil}
\externaldocument[adequate-]{adequate}
\externaldocument[dualizing-]{dualizing}
\externaldocument[duality-]{duality}
\externaldocument[discriminant-]{discriminant}
\externaldocument[derham-]{derham}
\externaldocument[local-cohomology-]{local-cohomology}
\externaldocument[algebraization-]{algebraization}
\externaldocument[curves-]{curves}
\externaldocument[resolve-]{resolve}
\externaldocument[models-]{models}
\externaldocument[functors-]{functors}
\externaldocument[equiv-]{equiv}
\externaldocument[pione-]{pione}
\externaldocument[etale-cohomology-]{etale-cohomology}
\externaldocument[proetale-]{proetale}
\externaldocument[relative-cycles-]{relative-cycles}
\externaldocument[more-etale-]{more-etale}
\externaldocument[trace-]{trace}
\externaldocument[crystalline-]{crystalline}
\externaldocument[spaces-]{spaces}
\externaldocument[spaces-properties-]{spaces-properties}
\externaldocument[spaces-morphisms-]{spaces-morphisms}
\externaldocument[decent-spaces-]{decent-spaces}
\externaldocument[spaces-cohomology-]{spaces-cohomology}
\externaldocument[spaces-limits-]{spaces-limits}
\externaldocument[spaces-divisors-]{spaces-divisors}
\externaldocument[spaces-over-fields-]{spaces-over-fields}
\externaldocument[spaces-topologies-]{spaces-topologies}
\externaldocument[spaces-descent-]{spaces-descent}
\externaldocument[spaces-perfect-]{spaces-perfect}
\externaldocument[spaces-more-morphisms-]{spaces-more-morphisms}
\externaldocument[spaces-flat-]{spaces-flat}
\externaldocument[spaces-groupoids-]{spaces-groupoids}
\externaldocument[spaces-more-groupoids-]{spaces-more-groupoids}
\externaldocument[bootstrap-]{bootstrap}
\externaldocument[spaces-pushouts-]{spaces-pushouts}
\externaldocument[spaces-chow-]{spaces-chow}
\externaldocument[groupoids-quotients-]{groupoids-quotients}
\externaldocument[spaces-more-cohomology-]{spaces-more-cohomology}
\externaldocument[spaces-simplicial-]{spaces-simplicial}
\externaldocument[spaces-duality-]{spaces-duality}
\externaldocument[formal-spaces-]{formal-spaces}
\externaldocument[restricted-]{restricted}
\externaldocument[spaces-resolve-]{spaces-resolve}
\externaldocument[formal-defos-]{formal-defos}
\externaldocument[defos-]{defos}
\externaldocument[cotangent-]{cotangent}
\externaldocument[examples-defos-]{examples-defos}
\externaldocument[algebraic-]{algebraic}
\externaldocument[examples-stacks-]{examples-stacks}
\externaldocument[stacks-sheaves-]{stacks-sheaves}
\externaldocument[criteria-]{criteria}
\externaldocument[artin-]{artin}
\externaldocument[quot-]{quot}
\externaldocument[stacks-properties-]{stacks-properties}
\externaldocument[stacks-morphisms-]{stacks-morphisms}
\externaldocument[stacks-limits-]{stacks-limits}
\externaldocument[stacks-cohomology-]{stacks-cohomology}
\externaldocument[stacks-perfect-]{stacks-perfect}
\externaldocument[stacks-introduction-]{stacks-introduction}
\externaldocument[stacks-more-morphisms-]{stacks-more-morphisms}
\externaldocument[stacks-geometry-]{stacks-geometry}
\externaldocument[moduli-]{moduli}
\externaldocument[moduli-curves-]{moduli-curves}
\externaldocument[examples-]{examples}
\externaldocument[exercises-]{exercises}
\externaldocument[guide-]{guide}
\externaldocument[desirables-]{desirables}
\externaldocument[coding-]{coding}
\externaldocument[obsolete-]{obsolete}
\externaldocument[fdl-]{fdl}
\externaldocument[index-]{index}

% Theorem environments.
%
\theoremstyle{plain}
\newtheorem{theorem}[subsection]{Theorem}
\newtheorem{proposition}[subsection]{Proposition}
\newtheorem{lemma}[subsection]{Lemma}

\theoremstyle{definition}
\newtheorem{definition}[subsection]{Definition}
\newtheorem{example}[subsection]{Example}
\newtheorem{exercise}[subsection]{Exercise}
\newtheorem{situation}[subsection]{Situation}

\theoremstyle{remark}
\newtheorem{remark}[subsection]{Remark}
\newtheorem{remarks}[subsection]{Remarks}

\numberwithin{equation}{subsection}

% Macros
%
\def\lim{\mathop{\mathrm{lim}}\nolimits}
\def\colim{\mathop{\mathrm{colim}}\nolimits}
\def\Spec{\mathop{\mathrm{Spec}}}
\def\Hom{\mathop{\mathrm{Hom}}\nolimits}
\def\Ext{\mathop{\mathrm{Ext}}\nolimits}
\def\SheafHom{\mathop{\mathcal{H}\!\mathit{om}}\nolimits}
\def\SheafExt{\mathop{\mathcal{E}\!\mathit{xt}}\nolimits}
\def\Sch{\mathit{Sch}}
\def\Mor{\mathop{\mathrm{Mor}}\nolimits}
\def\Ob{\mathop{\mathrm{Ob}}\nolimits}
\def\Sh{\mathop{\mathit{Sh}}\nolimits}
\def\NL{\mathop{N\!L}\nolimits}
\def\CH{\mathop{\mathrm{CH}}\nolimits}
\def\proetale{{pro\text{-}\acute{e}tale}}
\def\etale{{\acute{e}tale}}
\def\QCoh{\mathit{QCoh}}
\def\Ker{\mathop{\mathrm{Ker}}}
\def\Im{\mathop{\mathrm{Im}}}
\def\Coker{\mathop{\mathrm{Coker}}}
\def\Coim{\mathop{\mathrm{Coim}}}

% Boxtimes
%
\DeclareMathSymbol{\boxtimes}{\mathbin}{AMSa}{"02}

%
% Macros for moduli stacks/spaces
%
\def\QCohstack{\mathcal{QC}\!\mathit{oh}}
\def\Cohstack{\mathcal{C}\!\mathit{oh}}
\def\Spacesstack{\mathcal{S}\!\mathit{paces}}
\def\Quotfunctor{\mathrm{Quot}}
\def\Hilbfunctor{\mathrm{Hilb}}
\def\Curvesstack{\mathcal{C}\!\mathit{urves}}
\def\Polarizedstack{\mathcal{P}\!\mathit{olarized}}
\def\Complexesstack{\mathcal{C}\!\mathit{omplexes}}
% \Pic is the operator that assigns to X its picard group, usage \Pic(X)
% \Picardstack_{X/B} denotes the Picard stack of X over B
% \Picardfunctor_{X/B} denotes the Picard functor of X over B
\def\Pic{\mathop{\mathrm{Pic}}\nolimits}
\def\Picardstack{\mathcal{P}\!\mathit{ic}}
\def\Picardfunctor{\mathrm{Pic}}
\def\Deformationcategory{\mathcal{D}\!\mathit{ef}}
\title{A morphism of triangles whose first two components are homotopy equivalences has the third one too}
\author{The Stacks Project Authors}
\date{stacks-project:dga.tex:local-snapshot-20261012; GNU FDL 1.2 or later}
\begin{document}
\maketitle
\noindent\textit{Editorial note.} The counted statement is the author's own lemma environment labelled lemma-third-isomorphism (source lines 1015-1028); the proof environment begins two lines below it and is its own proof, closing at source line 1062.\par\medskip
\noindent\textit{Editorial note.} The proof invokes two author results of the same chapter by label, namely the functoriality of the cone and the nilpotence of a composite of maps of admissible sequences; both are retained as context spans, and the generator expands that context transitively, so every reference resolves inside the assembled package.\par\medskip
\noindent\textit{Editorial note.} The proof displays one Xy-pic diagram of three rows; it is part of the retained source range and is reproduced as written, so no figure file is needed.\par\medskip
\noindent\textit{Editorial note.} The retained preamble is the chapter's own minimal preamble (sources/stacks-dga/preamble.tex, 217 lines), a verbatim excerpt of the native Xy-pic setup, the native theorem declarations with the mathematical macros, and the project's own no-op reference environments; no macro was invented or redefined.\par\medskip
\noindent\textit{Editorial scope.} One original result of the Stacks Project's Differential Graded Algebra chapter, counted once: in a morphism of triangles of dg modules, if the first two maps are homotopy equivalences then so is the third.\par\medskip
\noindent\textit{Allowed standard prerequisites (editorial).}\begin{itemize}
\item Cone of a morphism and its triangle. The cone C(f) of a map of dg modules with its inclusion, projection and the triangle it forms.. The triangles and their cones are the objects of the morphism (a, b, c) (source lines 1019-1026).
\item Homotopy equivalence and inverse in the homotopy category. A map is a homotopy equivalence when it has an inverse in the homotopy category of dg modules.. The hypothesis gives inverses a\textasciicircum{}\{-1\} and b\textasciicircum{}\{-1\}, and the conclusion is that c has one (source lines 1027, 1031-1034, 1060-1061).
\item Functoriality of the cone and nilpotence of a composite. A homotopy commutative square induces a map of the cones, and a composite of two compatible maps of admissible sequences is homotopic to zero (the author's lemmas lemma-functorial-cone and lemma-nilpotent, retained as context).. The first builds c'; the second shows that the composites of c and c' have square zero, which is what makes 2 - c an inverse (source lines 1035, 1056-1061).
\end{itemize}\medskip
\begin{lemma}
\label{lemma-functorial-cone}
Let $(A, \text{d})$ be a differential graded algebra.
Suppose that
$$
\xymatrix{
K_1 \ar[r]_{f_1} \ar[d]_a & L_1 \ar[d]^b \\
K_2 \ar[r]^{f_2} & L_2
}
$$
is a diagram of homomorphisms of differential graded $A$-modules which is
commutative up to homotopy.
Then there exists a morphism $c : C(f_1) \to C(f_2)$ which gives rise to
a morphism of triangles
$$
(a, b, c) : (K_1, L_1, C(f_1), f_1, i_1, p_1) \to
(K_1, L_1, C(f_1), f_2, i_2, p_2)
$$
in $K(\text{Mod}_{(A, \text{d})})$.
\end{lemma}

\begin{proof}
Let $h : K_1 \to L_2$ be a homotopy between $f_2 \circ a$ and $b \circ f_1$.
Define $c$ by the matrix
$$
c =
\left(
\begin{matrix}
b & h \\
0 & a
\end{matrix}
\right) :
L_1 \oplus K_1 \to L_2 \oplus K_2
$$
A matrix computation show that $c$ is a morphism of differential
graded modules. It is trivial that $c \circ i_1 = i_2 \circ b$, and it is
trivial also to check that $p_2 \circ c = a \circ p_1$.
\end{proof}

\begin{lemma}
\label{lemma-nilpotent}
Let $(A, \text{d})$ be a differential graded algebra.
Let $0 \to K_i \to L_i \to M_i \to 0$, $i = 1, 2, 3$
be admissible short exact sequence of differential graded $A$-modules.
Let $b : L_1 \to L_2$ and $b' : L_2 \to L_3$
be homomorphisms of differential graded modules such that
$$
\vcenter{
\xymatrix{
K_1 \ar[d]_0 \ar[r] &
L_1 \ar[r] \ar[d]_b &
M_1 \ar[d]_0 \\
K_2 \ar[r] & L_2 \ar[r] & M_2
}
}
\quad\text{and}\quad
\vcenter{
\xymatrix{
K_2 \ar[d]^0 \ar[r] &
L_2 \ar[r] \ar[d]^{b'} &
M_2 \ar[d]^0 \\
K_3 \ar[r] & L_3 \ar[r] & M_3
}
}
$$
commute up to homotopy. Then $b' \circ b$ is homotopic to $0$.
\end{lemma}

\begin{proof}
By Lemma \ref{lemma-make-commute-map} we can replace $b$ and $b'$ by
homotopic maps such that the right square of the left diagram commutes
and the left square of the right diagram commutes. In other words, we have
$\Im(b) \subset \Im(K_2 \to L_2)$ and
$\Ker(b') \supset \Im(K_2 \to L_2)$.
Then $b' \circ b = 0$ as a map of modules.
\end{proof}

\begin{lemma}
\label{lemma-make-commute-map}
Let $(A, \text{d})$ be a differential graded algebra. Let
$$
\xymatrix{
K \ar[r]_f \ar[d]_a & L \ar[d]^b \\
M \ar[r]^g & N
}
$$
be a diagram of homomorphisms of differential graded $A$-modules
commuting up to homotopy.
\begin{enumerate}
\item If $f$ is an admissible monomorphism, then $b$ is homotopic to a
homomorphism which makes the diagram commute.
\item If $g$ is an admissible epimorphism, then $a$ is homotopic to a
morphism which makes the diagram commute.
\end{enumerate}
\end{lemma}

\begin{proof}
Let $h : K \to N$ be a homotopy between $bf$ and $ga$, i.e.,
$bf - ga = \text{d}h + h\text{d}$. Suppose that $\pi : L \to K$
is a graded $A$-module map left inverse to $f$. Take
$b' = b - \text{d}h\pi - h\pi \text{d}$.
Suppose $s : N \to M$ is a graded $A$-module map right inverse to $g$.
Take $a' = a + \text{d}sh + sh\text{d}$.
Computations omitted.
\end{proof}

\begin{lemma}
\label{lemma-third-isomorphism}
Let $(A, \text{d})$ be a differential graded algebra.
Let $f_1 : K_1 \to L_1$ and $f_2 : K_2 \to L_2$ be homomorphisms of
differential graded $A$-modules. Let
$$
(a, b, c) :
(K_1, L_1, C(f_1), f_1, i_1, p_1)
\longrightarrow
(K_1, L_1, C(f_1), f_2, i_2, p_2)
$$
be any morphism of triangles of $K(\text{Mod}_{(A, \text{d})})$.
If $a$ and $b$ are homotopy equivalences then so is $c$.
\end{lemma}

\begin{proof}
Let $a^{-1} : K_2 \to K_1$ be a homomorphism of differential graded $A$-modules
which is inverse to $a$ in $K(\text{Mod}_{(A, \text{d})})$.
Let $b^{-1} : L_2 \to L_1$ be a homomorphism of differential graded $A$-modules
which is inverse to $b$ in $K(\text{Mod}_{(A, \text{d})})$.
Let $c' : C(f_2) \to C(f_1)$ be the morphism from
Lemma \ref{lemma-functorial-cone} applied to
$f_1 \circ a^{-1} = b^{-1} \circ f_2$.
If we can show that $c \circ c'$ and $c' \circ c$ are isomorphisms in
$K(\text{Mod}_{(A, \text{d})})$
then we win. Hence it suffices to prove the following: Given
a morphism of triangles
$(1, 1, c) : (K, L, C(f), f, i, p)$
in $K(\text{Mod}_{(A, \text{d})})$ the morphism $c$ is an isomorphism
in $K(\text{Mod}_{(A, \text{d})})$.
By assumption the two squares in the diagram
$$
\xymatrix{
L \ar[r] \ar[d]_1 &
C(f) \ar[r] \ar[d]_c &
K[1] \ar[d]_1 \\
L \ar[r] &
C(f) \ar[r] &
K[1]
}
$$
commute up to homotopy. By construction of $C(f)$ the rows
form admissible short exact sequences. Thus we see that
$(c - 1)^2 = 0$ in $K(\text{Mod}_{(A, \text{d})})$ by
Lemma \ref{lemma-nilpotent}.
Hence $c$ is an isomorphism in $K(\text{Mod}_{(A, \text{d})})$
with inverse $2 - c$.
\end{proof}

\end{document}
